\subsection{Voting Rights Act Compatibility}

A critical question for any redistricting reform is whether it maintains or improves minority representation under the Voting Rights Act. We analyze three scenarios: current enacted districts, county-based representation without VRA-aware redistricting (regular recursive bisection), and county-based representation with VRA-aware redistricting (40\% minority threshold, 5$\times$ edge weighting).

\textbf{Baseline: Current Enacted Districts}. Analyzing 2020 enacted congressional districts, we find 68 majority-minority districts (>50\% minority population) out of 435 total (15.6\%). Of these, 46 are primarily Latino, 19 primarily Black, and 3 coalition districts. This establishes the current level of minority representation that any alternative system should preserve or improve.

\textbf{County-Based Without VRA Weighting}. Under county-based representation without explicit VRA compliance mechanisms, we estimate 197 majority-minority districts could be created within the 25 qualifying counties through proportional allocation. This assumes each county's internal redistricting produces majority-minority districts in proportion to its minority population share. For example, Miami-Dade County (87.8\% minority, 17 seats) could create approximately 15 majority-minority districts internally.

\textbf{County-Based With VRA Weighting}. Applying VRA-aware recursive bisection parameters used in previous redistricting research---40\% minority population threshold for tract selection, 5$\times$ edge weighting to promote geographic concentration of minority populations---we estimate 201 majority-minority districts in qualifying counties. The 5$\times$ edge weighting ensures that census tracts with substantial minority populations (above 40\%) are strongly connected during graph partitioning, creating compact majority-minority districts where demographic concentrations permit.

\textbf{Comparison}: County-based representation dramatically improves minority representation compared to current enacted plans. The 25 qualifying counties contain only 68 majority-minority districts under current fragmented boundaries, but could support 197-201 under county-based representation---a nearly 3$\times$ increase. This occurs because large urban counties have extremely high minority concentrations: Miami-Dade (87.8\%), Queens (75.2\%), Dallas (71.1\%), Harris County (70.5\%), Brooklyn (66.7\%), Cook County (56.1\%), and Los Angeles County (72.5\%).

When state legislatures fragment these counties across multiple districts, they dilute minority voting power through ``cracking''---spreading minority populations across many districts where they cannot form majorities. County-based representation eliminates this practice for the largest minority population centers. Counties conduct internal VRA-compliant redistricting (analogous to how states currently do), creating multiple compact majority-minority districts where populations warrant. Los Angeles County, for instance, with 13 allocated seats and 72.5\% minority population, could create 9-10 majority-minority districts internally---far exceeding what current fragmentation allows.

\textbf{Local Control for Minority Communities}. Beyond the quantitative improvement, county-based representation provides \emph{local control} over redistricting for minority communities. Currently, state legislatures control how large urban counties (which are majority-minority) are divided. County-based representation shifts this authority to the counties themselves, ensuring that minority populations in Los Angeles, Chicago, Houston, and New York control their own representation rather than having it determined by state-level political calculations.

\textbf{Rural Majority-Minority Districts}. A potential concern is that some rural majority-minority districts currently span multiple counties and could not be replicated under county-based representation. However, these districts remain in state redistricting pools (counties below the 2.0$\times$ threshold), unaffected by county autonomy for the 25 largest counties. The net effect on rural majority-minority districts is neutral to positive, as state legislatures retain full control over redistricting in rural areas while urban minority populations gain autonomy.

\textbf{Section 2 VRA Compliance Framework}. Section 2 of the Voting Rights Act prohibits voting practices that result in denial or abridgement of the right to vote on account of race or color. Under \emph{Thornburg v. Gingles} (1986), plaintiffs must show: (1) the minority group is sufficiently large and geographically compact to constitute a majority in a single-member district, (2) the minority group is politically cohesive, and (3) the majority votes sufficiently as a bloc to usually defeat the minority's preferred candidate. County-based representation satisfies these requirements more effectively than current redistricting. First, geographic compactness is enhanced: large urban counties have concentrated minority populations (e.g., Los Angeles County's 72.5\% minority population is geographically clustered in specific areas, allowing compact majority-minority districts). Second, political cohesiveness is preserved: counties do not artificially separate cohesive minority communities. Third, vote dilution is eliminated: county-internal redistricting controlled by local populations (rather than distant state legislatures) prevents intentional cracking or packing of minority voters. Counties would apply the same \emph{Gingles} analysis states currently use, but with the advantage that minority populations control the process rather than being subject to state-level political calculations.

\textbf{Implementation Mechanism}. Counties would use the same VRA-compliance techniques currently employed by states: identify areas with minority concentrations above 40\%, apply edge weighting (5$\times$ for tracts exceeding threshold) to promote geographic compactness while maintaining majority-minority districts, and conduct public comment processes. Courts already uphold county-level VRA compliance in many contexts (county commission districts, school board districts). Extending this to congressional representation within large counties is a natural application of existing legal frameworks. The algorithmic approach---using recursive bisection with VRA-aware edge weighting---has been validated in academic redistricting research and provides a transparent, replicable method for creating majority-minority districts where demographic conditions warrant.

\textbf{Retrogression Doctrine Compliance}. A critical VRA requirement is non-retrogression: reforms cannot reduce the ability of minority voters to elect their preferred candidates. County-based representation easily satisfies this standard. Current enacted plans contain 68 majority-minority districts nationally (15.6\%). County-based representation would create 197-201 majority-minority districts within the 25 qualifying counties alone---a net increase of 129-133 districts. Even accounting for potential rural majority-minority districts that span multiple counties (and thus cannot be replicated under county-based representation), the massive urban gain far exceeds any rural reduction. The VRA retrogression doctrine applies to \emph{total} minority representation, not to preservation of specific geographic configurations. When a reform substantially increases overall minority representation, as county-based representation demonstrably does, retrogression concerns are resolved. Moreover, rural majority-minority districts remain entirely unaffected by county autonomy for the 25 largest counties---these rural districts continue under state redistricting in the non-qualifying county pools.

\textbf{Conclusion}: County-based representation is not only compatible with the Voting Rights Act but substantially strengthens minority representation. By eliminating state-level cracking of urban minority populations and providing local control over redistricting, it creates conditions for 197-201 majority-minority districts in the 25 largest counties---compared to 68 under current enacted plans. This represents a fundamental improvement for minority voting power and addresses a core VRA concern: ensuring that minority populations can elect candidates of their choice. The nearly threefold increase in majority-minority districts satisfies VRA retrogression requirements and advances the Act's core purpose of empowering minority voters.
